A fazer
-------
- Padronizar hebraico em alfabeto latino e também a ordem das palavras traduzidas, como "em português…";
- Padronizar caixa alta e baixa (bibliografia; conceitos como Tempo, Espaço; etc);
- Checar primeiras citações de bibliografia (apud; op. cit.; etc);
- Padronizar notas de rodapé, que devem ser indicadas após pontuação;
- Shabat e outras palavras serão italizadas? Em textos específicos como este, podemos escolher não italizar (já que são muitas as menções);
- Eu colocaria mais seções no texto. Ou dividiria em mais capítulos, mais explicativos. Eles poderiam até levar nomes mais específicos, que dizem respeito ao que estão falando;
- Midrash é um livro religioso? E a Mishná e a Guemará? Os outros (Talmud, Torá) foram desitalizados, conforme a tradição tipográfica.
- Cortar e ajeitar texto.
- Falar na introdução sobre o tipo de escrita fragmentária. E assumir isso no projeto gráfico, contar com ele pra realçar e dar vida ao texto.
- Pensar em jogar as notas maiores pro final. Manter apenas as pequenas perto do texto. Também podemos fazer uma edição das notas. Ou fazer um glossário, simplesmente.
- Pensar no que compor o apêndice.
- Escrever melhor as didascálias.

Já feito
--------
- Emendas da preparadora batidas. Ela fez um bom trabalho, mas se manteve no básico;
- "Tradução nossa" virou "Tradução da autora";
- Ajustei as notas de rodapé, que estavam com uma numeração estranh, tipo "3 1/2". Algumas não estavam linkadas, fiz à mão. Então precisa de revisão;
- Retirei todos os itálicos que vinham com aspas.